\section{When the financial history itself cannot be settled}
\label{sim:financing-viability-boundary}

The receipt-6 contradiction contains an economic impossibility that additional information cannot remove. Its two second-stage refusal rows share the exact rate point $19/320$ and $43/800$. Under the existing floating-coupon formula, these rates give $I_1=95/16$ and $I_2=43/8$. Checking all four action sequences at this one point avoids confusing two different financial histories that happen to produce the same message.

Opening cash $11/2=88/16$ cannot pay the first coupon $95/16$, so either sequence beginning with pay is refused immediately. Drawing supplies one unit and costs $1/4$, leaving first-settlement cash $5/16$. Receipt 6 followed by repayment of the borrowed principal leaves $85/16$, while the second coupon is $86/16$. Draw/pay therefore fails by $1/16$. Draw/draw also fails: the once-only facility has already been used, and the second coupon is below its strict draw trigger. All four sequences fail even for a manager who knows both coupons in advance.

At the same rate point with receipt 7, draw/pay succeeds and ends with cash $15/16$. This is an explicit counterfactual, not a replacement for the assumed receipt-6 world. The earlier viable receipt-6 examples also remain valid. The diagnostic establishes one impossible point in the original family, rather than impossibility at every point. Later algebraic cash from a sequence whose first payment was refused is not an executed outcome.\footnote{\path{data/accounting-master/policy-failure-diagnosis-20261027/} retains ten exact path comparisons, 337 assertions and seven alteration controls. Complete states and all guards agree with the independent recurrence; the two second-stage refusal rows use the same rational witness.}

\subsection{An exact boundary on the committed-right coupon domain}

Keep the committed right, opening cash, once-only borrowing rule, fee and repayment timing unchanged. Let both coupons range over $[3,6]$, and take $R$ to be either 6 or 7. If $I_1\le 11/2$, the first payment can be made without borrowing. Cash after the receipt is then $11/2-I_1+R\ge R\ge6$, enough to pay every permitted second coupon. Thus pay/pay is viable throughout this first region.

If $I_1>11/2$, paying first is impossible, while the trigger permits a first draw. Since $I_1\le6$, the net draw of $3/4$ leaves nonnegative cash after that payment. Receipt and repayment then leave
\begin{equation}
 C_{\mathrm{before\ second}}=\frac{11}{2}-I_1+R-\frac14
 =R+\frac{21}{4}-I_1.
 \label{eq:viability-second-cash}
\end{equation}
The first borrowing has used the facility, so the second coupon must be paid from this cash. Draw/pay is therefore viable exactly when $I_1+I_2\le R+21/4$. Combining the two regions, the economically infeasible set is
\begin{equation}
 \mathcal N_R=\left\{(I_1,I_2)\in[3,6]^2:
 I_1>\frac{11}{2},\quad I_1+I_2>R+\frac{21}{4}\right\}.
 \label{eq:committed-infeasible-region}
\end{equation}
For receipt 6, the sum boundary is $45/4$. Equality is viable: the second payment leaves exactly zero cash. Equality at the first trigger is also viable, because $I_1=11/2$ can be paid without drawing. For receipt 7, the sum boundary is $49/4>12$, above the largest possible coupon sum, so $\mathcal N_7$ is empty on this domain.

The result concerns the continuous coupon square. Exact affine certificates check every inherited guard of pay/pay in the first region and draw/pay in the second, and exclude every action sequence in $\mathcal N_R$. The three regions are exhaustive and disjoint because they split first on $I_1\le11/2$ and then on the sum inequality. Rational boundary examples illustrate the result, while the region certificates establish it.\footnote{\path{data/accounting-master/financing-viability-boundary-20261027/} checks two receipt domains, 366 guard implications, eight action-region exclusions and eight boundary examples, with 315 assertions and ten alteration controls. Emptiness certificates are verified independently of their producer.}

On the viable subset, the original refined first trigger suffices for a common rule: draw first exactly when that trigger is true, then pay second. The manager does not need the future coupon to choose the first action. This rule cannot settle the excluded points, which remain part of the original model. The result does not cover absent or mixed borrowing rights, arbitrary receipts, or a changed facility contract.

A financially impossible settlement is still a possible company history to represent. The unpaid coupon, the earlier draw and repayment, and the dates of failed payments must remain available to subsequent accounting and recovery events. Refusing a payment instruction must therefore preserve the outstanding obligation and the failed history rather than remove it from later statements.
